% ============================================================
%  Experiment 2 --- the reprojection Jacobian and pose observability
%  Standalone note. Compiles with pdfLaTeX (on Overleaf: just upload it).
%  Study aid kept alongside the code; it is not included in report.tex.
% ============================================================
\documentclass[11pt,a4paper]{article}

\usepackage[utf8]{inputenc}
\usepackage[T1]{fontenc}
\usepackage[english]{babel}
\usepackage[a4paper,margin=2.5cm]{geometry}
\usepackage{amsmath,amssymb,amsthm,mathtools}
\usepackage{booktabs,array,longtable}
\usepackage{enumitem}
\usepackage[colorlinks=true,linkcolor=blue!55!black]{hyperref}

\theoremstyle{remark}
\newtheorem*{remark}{Remark}

% --- notation shortcuts (none of these clash with LaTeX built-ins) ---
\newcommand{\Rset}{\mathbb{R}}
\newcommand{\Rmat}{\mathbf{R}}
\newcommand{\Jmat}{\mathbf{J}}
\newcommand{\Kmat}{\mathbf{K}}
\newcommand{\Dmat}{\mathbf{D}}
\newcommand{\Umat}{\mathbf{U}}
\newcommand{\Vmat}{\mathbf{V}}
\newcommand{\Smat}{\mathbf{S}}
\newcommand{\Mmat}{\mathbf{M}}
\newcommand{\Sigmat}{\boldsymbol{\Sigma}}
\newcommand{\Idm}{\mathbf{I}}
\newcommand{\evec}{\mathbf{e}}
\newcommand{\nvec}{\mathbf{n}}
\newcommand{\Pvec}{\mathbf{P}}
\newcommand{\pvec}{\mathbf{p}}
\newcommand{\tvec}{\mathbf{t}}
\newcommand{\uvec}{\mathbf{u}}
\newcommand{\vvec}{\mathbf{v}}
\newcommand{\wvec}{\mathbf{w}}
\newcommand{\Xvec}{\mathbf{X}}
\newcommand{\xvec}{\mathbf{x}}
\newcommand{\bvec}{\mathbf{b}}
\newcommand{\avec}{\mathbf{a}}
\newcommand{\zerovec}{\mathbf{0}}
\newcommand{\thetavec}{\boldsymbol{\theta}}
\newcommand{\dth}{\delta\boldsymbol{\theta}}
\newcommand{\omegavec}{\boldsymbol{\omega}}
\newcommand{\dom}{\delta\boldsymbol{\omega}}
\newcommand{\dtv}{\delta\mathbf{t}}
\newcommand{\skewm}[1]{[#1]_{\times}}
\newcommand{\Cov}{\operatorname{Cov}}

\title{Experiment 2: the Reprojection Jacobian\\[0.3em]
       \large Derivation, singular values, and what they say about observability}
\date{}

\begin{document}
\maketitle

\noindent
The assignment frames pose estimation from an ArUco marker as the minimization
of the reprojection error
\[
  E(\Rmat,\tvec) = \sum_i \big\| \xvec_i - \pi\big(\Kmat(\Rmat\Xvec_i + \tvec)\big)\big\|^2 ,
\]
and asks, for every viewing angle, for the Jacobian of the residuals with
respect to the six pose parameters $\thetavec = (t_x,t_y,t_z,\omega_x,\omega_y,
\omega_z)$, its singular values, its smallest singular value and its condition
number, in order to find the directions of the pose that are poorly observable.
This note derives all of it. Sections~\ref{sec:problem}--\ref{sec:units} are
general; Section~\ref{sec:frontal} solves the fronto-parallel case in closed
form; Section~\ref{sec:sweep} gives the numbers along the camera arc.

\medskip\noindent
\emph{Notation.} The assignment writes the residual vector as $\mathbf{r}$. We
call it $\evec$, for consistency with the report; nothing else changes.

\tableofcontents

% ============================================================
\section{Symbols}
\label{sec:symbols}
% ============================================================

\begin{longtable}{@{}lp{11.2cm}@{}}
\toprule
Symbol & Meaning \\
\midrule
\endfirsthead
\toprule
Symbol & Meaning \\
\midrule
\endhead
$N$ & number of corners; here $N = 4$ \\
$\Xvec_i$ & corner $i$ in the marker frame, $\Xvec_i = (X_i,Y_i,0)^{\top}$ \\
$\pvec_i$ & observed pixel of corner $i$ (the $\xvec_i$ of the assignment) \\
$\Rmat,\ \tvec$ & pose, marker frame $\to$ camera frame \\
$\Pvec_i$ & corner $i$ in the camera frame, $\Pvec_i = \Rmat\Xvec_i + \tvec
  = (P_{i,1},P_{i,2},P_{i,3})^{\top}$ \\
$\Kmat$ & intrinsics, with focal lengths $f_x,f_y$ and principal point
  $(c_x,c_y)$ \\
$\pi$ & perspective division, $\pi(a,b,c) = (a/c,\ b/c)$ \\
$\evec_i$ & residual of corner $i$, a vector in $\Rset^2$ (pixels) \\
$\evec$ & all residuals stacked, a vector in $\Rset^{2N} = \Rset^8$ \\
$E$ & reprojection cost, $E = \|\evec\|^2$ (pixels$^2$) \\
$\omegavec$ & local rotation parameters (radians); see
  Section~\ref{sec:param} \\
$\thetavec$ & the six pose parameters, $\thetavec = (\tvec,\omegavec)$ \\
$\dth = (\dtv,\dom)$ & a small change of the six parameters \\
$\skewm{\avec}$ & skew-symmetric matrix with $\skewm{\avec}\wvec =
  \avec\times\wvec$ \\
$\exp(\cdot)$ & matrix exponential; $\exp(\skewm{\omegavec})$ is the rotation by
  angle $\|\omegavec\|$ about the axis $\omegavec/\|\omegavec\|$ \\
$\Dmat_i$ & $2\times3$ derivative of the pixel projection with respect to
  $\Pvec_i$ \\
$\Jmat$ & the $8\times6$ Jacobian $\partial\evec/\partial\thetavec$ \\
$\Jmat_i$ & the two rows of $\Jmat$ belonging to corner $i$ \\
$\Umat,\Sigmat,\Vmat$ & singular value decomposition $\Jmat =
  \Umat\Sigmat\Vmat^{\top}$ \\
$\sigma_k$ & $k$-th singular value of $\Jmat$,
  $\sigma_1 \ge \dots \ge \sigma_6 \ge 0$ \\
$\vvec_k$ & $k$-th right singular vector: a direction in the space of the six
  pose parameters \\
$\uvec_k$ & $k$-th left singular vector: a direction in the space of the eight
  residuals \\
$\kappa$ & condition number, $\kappa = \sigma_1/\sigma_6$; from
  Section~\ref{sec:units} on, of the dimensionless Jacobian $\Jmat_{\mathrm{n}}$ \\
$\Jmat_{\mathrm{n}}$ & dimensionless Jacobian, $\Jmat_{\mathrm{n}} = \Jmat\Smat$,
  with translation measured in units of $Z$ \\
$\kappa_{\mathrm{SI}}$ & condition number of $\Jmat$ itself, translation in metres
  and rotation in radians \\
$\nvec$ & pixel noise on the eight observed coordinates \\
$\sigma_{\mathrm{px}}$ & standard deviation of each noise component (pixels) \\
$\Smat$ & diagonal rescaling of the parameters,
  $\Smat = \operatorname{diag}(Z,Z,Z,1,1,1)$; see Section~\ref{sec:units} \\
$\theta_v$ & viewing angle (to avoid a clash with $\thetavec$) \\
$h$ & half the side of the square, $h = L/2$ \\
$Z$ & distance from the camera to the marker centre \\
$f$ & focal length when $f_x = f_y$ \\
$s$ & the length $s = h^2/Z$, which appears in Section~\ref{sec:frontal} \\
\bottomrule
\end{longtable}

% ============================================================
\section{The estimation problem}
\label{sec:problem}
% ============================================================

For a pose $(\Rmat,\tvec)$, corner $i$ lands in the camera frame at
$\Pvec_i = \Rmat\Xvec_i + \tvec$ and is projected to the pixel
\begin{equation}
  \hat\pvec_i
  = \pi\big(\Kmat\Pvec_i\big)
  = \begin{pmatrix}
      f_x\,P_{i,1}/P_{i,3} + c_x \\[2pt]
      f_y\,P_{i,2}/P_{i,3} + c_y
    \end{pmatrix}.
  \label{eq:proj}
\end{equation}
Its residual is the difference between prediction and observation,
\begin{equation}
  \evec_i = \hat\pvec_i - \pvec_i \in \Rset^2 ,
  \label{eq:residual}
\end{equation}
and stacking the four corners gives
\begin{equation}
  \evec = \begin{pmatrix} \evec_1 \\ \evec_2 \\ \evec_3 \\ \evec_4 \end{pmatrix}
  \in \Rset^8,
  \qquad
  E(\Rmat,\tvec) = \sum_{i=1}^{4}\|\evec_i\|^2 = \|\evec\|^2 .
  \label{eq:cost}
\end{equation}
This is the cost of the assignment (the sign of $\evec_i$ does not matter,
since it is squared). For reference, the RMSE plotted in Experiment~1 is
$\sqrt{E/8}$.

\paragraph{Counting.} Four corners give eight equations; the pose has six
unknowns. The problem is therefore over-determined by two, which is why a
least-squares formulation is needed at all, and why noise cannot in general be
fitted exactly.

% ============================================================
\section{Parametrizing the pose locally}
\label{sec:param}
% ============================================================

Translation is a vector, and a small change is simply
$\tvec \mapsto \tvec + \dtv$. Rotation is not: the set of rotations is curved,
and adding a small matrix to $\Rmat$ leaves it. Euler angles would avoid that
but introduce singularities (gimbal lock) and make the Jacobian depend on an
arbitrary axis order. The standard remedy is a \emph{local} parametrization
around the current rotation:
\begin{equation}
  \Rmat(\dom) = \Rmat\,\exp\!\big(\skewm{\dom}\big),
  \qquad
  \skewm{\avec} =
  \begin{bmatrix}
    0 & -a_3 & a_2 \\
    a_3 & 0 & -a_1 \\
    -a_2 & a_1 & 0
  \end{bmatrix}.
  \label{eq:localrot}
\end{equation}
For $\dom = \zerovec$ this is $\Rmat$ itself; for small $\dom$ it is a rotation
by angle $\|\dom\|$ about the axis $\dom/\|\dom\|$, applied \emph{first}, in the
marker frame, because it multiplies $\Rmat$ on the right. The six parameters are
\begin{equation}
  \dth = (\delta t_x,\delta t_y,\delta t_z,\ \delta\omega_x,\delta\omega_y,
  \delta\omega_z)^{\top},
  \label{eq:theta}
\end{equation}
with translations in metres and rotations in radians. To first order the
exponential is the identity plus its argument,
\begin{equation}
  \exp\!\big(\skewm{\dom}\big) = \Idm + \skewm{\dom} + O(\|\dom\|^2),
  \label{eq:expfirst}
\end{equation}
which is all the Jacobian needs. This is the parametrization implemented in
\texttt{jacobian.py} (\texttt{perturb\_pose}).

% ============================================================
\section{Deriving the Jacobian}
\label{sec:jacobian}
% ============================================================

The residual depends on the pose only through $\Pvec_i$, so by the chain rule
\begin{equation}
  \Jmat_i = \frac{\partial \evec_i}{\partial \dth}
  = \underbrace{\frac{\partial \hat\pvec_i}{\partial \Pvec_i}}_{\Dmat_i\ (2\times3)}
    \;
    \underbrace{\frac{\partial \Pvec_i}{\partial \dth}}_{(3\times6)} .
  \label{eq:chain}
\end{equation}
The observation $\pvec_i$ is a constant and does not contribute.

\subsection{How the camera-frame point moves}

Apply the perturbation to $\Pvec_i$:
\begin{align}
  \Pvec_i(\dth)
  &= \Rmat\,\exp\!\big(\skewm{\dom}\big)\,\Xvec_i + \tvec + \dtv \nonumber \\
  &\overset{\eqref{eq:expfirst}}{=}
     \Rmat\Xvec_i + \tvec + \dtv + \Rmat\,\skewm{\dom}\,\Xvec_i
     + O(\|\dth\|^2) .
\end{align}
The last term is $\Rmat(\dom\times\Xvec_i)$. Since the cross product is
antisymmetric, $\dom\times\Xvec_i = -\Xvec_i\times\dom = -\skewm{\Xvec_i}\dom$,
so
\begin{equation}
  \Pvec_i(\dth) = \Pvec_i + \dtv - \Rmat\,\skewm{\Xvec_i}\,\dom
  + O(\|\dth\|^2),
\end{equation}
and reading off the coefficients of $\dtv$ and $\dom$,
\begin{equation}
  \boxed{\;
  \frac{\partial \Pvec_i}{\partial \dth}
  = \begin{bmatrix} \Idm_3 & -\Rmat\,\skewm{\Xvec_i} \end{bmatrix}
  \;}
  \qquad (3\times6).
  \label{eq:dPdtheta}
\end{equation}

\subsection{How the pixel moves with the camera-frame point}

Differentiate \eqref{eq:proj} with respect to $\Pvec_i$. With
$(X,Y,Z) = (P_{i,1},P_{i,2},P_{i,3})$ for brevity, the quotient rule gives
\begin{equation}
  \boxed{\;
  \Dmat_i = \frac{\partial \hat\pvec_i}{\partial \Pvec_i}
  = \begin{bmatrix}
      f_x/Z & 0 & -f_x X/Z^2 \\[2pt]
      0 & f_y/Z & -f_y Y/Z^2
    \end{bmatrix}
  \;}
  \qquad (2\times3).
  \label{eq:D}
\end{equation}
The first two columns say that moving the point sideways moves its pixel by
$f/Z$ per metre; the third says that moving it away from the camera pulls the
pixel towards the principal point.

\subsection{Assembling}

Combining \eqref{eq:chain}, \eqref{eq:dPdtheta} and \eqref{eq:D},
\begin{equation}
  \boxed{\;
  \Jmat_i = \Dmat_i
  \begin{bmatrix} \Idm_3 & -\Rmat\,\skewm{\Xvec_i} \end{bmatrix}
  \;}
  \qquad (2\times6),
  \label{eq:Ji}
\end{equation}
and stacking the four corners,
\begin{equation}
  \Jmat = \begin{bmatrix} \Jmat_1 \\ \Jmat_2 \\ \Jmat_3 \\ \Jmat_4 \end{bmatrix}
  \in \Rset^{8\times6}.
  \label{eq:J}
\end{equation}
So the assignment's guess is right: the Jacobian is $8\times6$. Its first three
columns have units of pixels per metre, its last three pixels per radian. The
implementation is \texttt{reprojection\_jacobian} in \texttt{jacobian.py}, and
the test suite checks it against centred finite differences to a relative
accuracy of $2\cdot10^{-6}$.

% ============================================================
\section{What the singular values say}
\label{sec:svd}
% ============================================================

\subsection{The decomposition}

Every real $8\times6$ matrix has a singular value decomposition
\begin{equation}
  \Jmat = \Umat\,\Sigmat\,\Vmat^{\top}
  = \sum_{k=1}^{6} \sigma_k\,\uvec_k\,\vvec_k^{\top},
  \label{eq:svd}
\end{equation}
with $\Vmat = [\vvec_1\cdots\vvec_6]$ orthonormal in the space of pose
parameters, $\Umat = [\uvec_1\cdots\uvec_6]$ orthonormal columns in the space of
residuals, and $\sigma_1 \ge \dots \ge \sigma_6 \ge 0$. Multiplying by
$\vvec_k$,
\begin{equation}
  \Jmat\,\vvec_k = \sigma_k\,\uvec_k
  \qquad\Longrightarrow\qquad
  \|\Jmat\,\vvec_k\| = \sigma_k .
  \label{eq:Jvk}
\end{equation}

\paragraph{Reading.} Move the pose by a small amount $\varepsilon$ along
$\vvec_k$. To first order the residuals change by $\varepsilon\,\sigma_k$ in
norm. A large $\sigma_k$ means that direction of the pose is clearly visible
in the image; a small one means the pose can move along $\vvec_k$ while the four
corners barely move. That is precisely a \emph{poorly observable direction}:
\begin{itemize}[leftmargin=*]
  \item $\sigma_6 = \sigma_{\min}$ measures how poorly the worst direction is
        observed;
  \item $\vvec_6$ says which combination of the six parameters that worst
        direction is;
  \item $\kappa = \sigma_1/\sigma_6$ compares the best and worst observed
        directions; a large $\kappa$ means the problem is badly balanced.
\end{itemize}
If $\sigma_6 = 0$ the pose would be locally unidentifiable along $\vvec_6$. For
four corners of a square in front of the camera this does not happen short of
the edge-on limit, so $\Jmat$ has full rank $6$.

\subsection{From singular values to uncertainty}

Why $\sigma_k$ is the right measure becomes precise once noise is added. Let
$\thetavec^\star$ be the true pose, so that $\evec(\thetavec^\star) =
\zerovec$ on exact data, and let the observed pixels carry noise $\nvec$. The
residuals become $\evec(\thetavec) = \evec_0(\thetavec) - \nvec$, where
$\evec_0$ is the noise-free residual. Linearizing around the truth,
\begin{equation}
  \evec(\thetavec^\star + \dth) \approx \Jmat\,\dth - \nvec ,
\end{equation}
and the minimizer of $E = \|\evec\|^2$ is the linear least-squares solution
\begin{equation}
  \dth = \big(\Jmat^{\top}\Jmat\big)^{-1}\Jmat^{\top}\nvec .
  \label{eq:gn}
\end{equation}
This is also exactly one Gauss--Newton step. If the noise components are
independent with variance $\sigma_{\mathrm{px}}^2$, that is
$\Cov(\nvec) = \sigma_{\mathrm{px}}^2\Idm_8$, then
\begin{equation}
  \Cov(\dth)
  = \big(\Jmat^{\top}\Jmat\big)^{-1}\Jmat^{\top}
    \big(\sigma_{\mathrm{px}}^2\Idm\big)
    \Jmat\big(\Jmat^{\top}\Jmat\big)^{-1}
  = \sigma_{\mathrm{px}}^2\big(\Jmat^{\top}\Jmat\big)^{-1}.
  \label{eq:cov}
\end{equation}
Substituting \eqref{eq:svd}, $\Jmat^{\top}\Jmat = \Vmat\Sigmat^2\Vmat^{\top}$,
so
\begin{equation}
  \boxed{\;
  \Cov(\dth) = \sigma_{\mathrm{px}}^2 \sum_{k=1}^{6}
    \frac{\vvec_k\vvec_k^{\top}}{\sigma_k^2}
  \;}
  \label{eq:covsvd}
\end{equation}
The estimate scatters along each $\vvec_k$ with standard deviation
$\sigma_{\mathrm{px}}/\sigma_k$. The weakest direction $\vvec_6$ carries the
largest scatter, $\sigma_{\mathrm{px}}/\sigma_{\min}$: small singular values are
not an abstraction, they are where the pose error goes.

\begin{remark}
Equation \eqref{eq:cov} describes the minimizer of $E$, linearized around the
truth. It is a statement about the least-squares estimator the assignment
defines, not about any solver. A Monte Carlo check must therefore re-estimate
the pose by minimizing $E$ (for instance Gauss--Newton started from IPPE); raw
IPPE is algebraic and need not scatter as \eqref{eq:cov} predicts.
\end{remark}

% ============================================================
\section{Units, and the convention adopted}
\label{sec:units}
% ============================================================

\subsection{The problem}

The six parameters do not share a unit: three are metres, three radians. The
singular values of $\Jmat$ therefore mix pixels per metre with pixels per
radian, and $\kappa$ depends on that choice. Rescaling the parameters by a
diagonal matrix $\Smat$ replaces $\Jmat$ by $\Jmat\Smat$, which has different
singular values and a different $\kappa$, although nothing physical has
changed. Two examples make the point.
\begin{itemize}[leftmargin=*]
  \item \emph{Changing units.} Measuring translation in millimetres instead of
        metres moves $\kappa_{\mathrm{SI}}$ at a $30^{\circ}$ view from $55$ to
        $652$.
  \item \emph{The same image.} A $10$~cm marker at $0.5$~m and a $20$~cm marker
        at $1$~m produce exactly the same picture, yet head-on they give
        $\kappa_{\mathrm{SI}} = 200$ and $100$ respectively.
\end{itemize}
A quantity that differs for two identical images is not measuring the
geometry. It also makes some comparisons meaningless: in metres the depth
singular value exceeds the spin one, in centimetres it falls below it.

\subsection{The convention}

We measure translation in units of the camera--marker distance
$Z = \|\tvec\|$:
\begin{equation}
  \Jmat_{\mathrm{n}} = \Jmat\,\Smat,
  \qquad
  \Smat = \operatorname{diag}(Z,Z,Z,1,1,1).
  \label{eq:Jn}
\end{equation}
A translation of one unit then displaces the marker by about one radian as seen
from the camera, so both halves of $\thetavec$ are dimensionless and directly
comparable. Consequences:
\begin{itemize}[leftmargin=*]
  \item the singular values of $\Jmat_{\mathrm{n}}$ are in \emph{pixels}, and
        $\kappa$ is a pure number;
  \item two setups with the same image get the same $\kappa$ (both examples
        above give $100$), and $\kappa$ does not depend on the focal length,
        which scales every singular value alike;
  \item $\kappa$ still depends on the geometry, as it should: it is set by how
        large the marker appears (Section~\ref{sec:frontal});
  \item the rotation columns are untouched, so the singular values and
        directions of pure rotations do not change.
\end{itemize}
From here on, $\sigma_k$, $\vvec_k$ and $\kappa$ refer to $\Jmat_{\mathrm{n}}$,
and $\kappa_{\mathrm{SI}}$ to $\Jmat$. This is what
\texttt{run\_conditioning\_sweep} in \texttt{experiment2.py} computes. The
uncertainty statement \eqref{eq:covsvd} is unaffected, since it rescales
consistently; the Monte Carlo therefore keeps working with $\Jmat$ in SI units,
because its covariances are in metres and radians.

% ============================================================
\section{The fronto-parallel case, in closed form}
\label{sec:frontal}
% ============================================================

At a fronto-parallel view every quantity can be computed by hand, and the
result explains the numbers of Section~\ref{sec:sweep}.

\subsection{Setup}

Take $f_x = f_y = f$, the marker centred on the optical axis at distance $Z$ and
facing the camera,
\begin{equation}
  \Rmat = \operatorname{diag}(1,-1,-1),
  \qquad
  \tvec = (0,0,Z)^{\top},
\end{equation}
and the corners $\Xvec_i = (x,y,0)^{\top}$ with
$(x,y) \in \{(-h,h),\,(h,h),\,(h,-h),\,(-h,-h)\}$. This is exactly the
$0^{\circ}$ pose of the camera arc. Then
\begin{equation}
  \Pvec_i = \Rmat\Xvec_i + \tvec = (x,\ -y,\ Z)^{\top}:
\end{equation}
all four corners are at the same depth $Z$.

\subsection{The two rows of each corner}

From \eqref{eq:D} with $P_{i,1} = x$, $P_{i,2} = -y$, $P_{i,3} = Z$,
\begin{equation}
  \Dmat_i = \frac{f}{Z}
  \begin{bmatrix} 1 & 0 & -x/Z \\ 0 & 1 & y/Z \end{bmatrix}.
\end{equation}
For the rotation block, $\skewm{\Xvec_i}$ with $\Xvec_i = (x,y,0)^{\top}$ and
$\Rmat = \operatorname{diag}(1,-1,-1)$ give
\begin{equation}
  \skewm{\Xvec_i} =
  \begin{bmatrix} 0 & 0 & y \\ 0 & 0 & -x \\ -y & x & 0 \end{bmatrix},
  \qquad
  -\Rmat\,\skewm{\Xvec_i} =
  \begin{bmatrix} 0 & 0 & -y \\ 0 & 0 & -x \\ -y & x & 0 \end{bmatrix}.
\end{equation}
Multiplying out \eqref{eq:Ji}, the two rows of corner $i$, in the parameter
order $(t_x,t_y,t_z,\omega_x,\omega_y,\omega_z)$, are
\begin{align}
  \text{$u$-row:}\quad
  &\frac{f}{Z}\Big(\,1,\ \ 0,\ \ -\frac{x}{Z},\ \ \frac{xy}{Z},\ \
     -\frac{x^2}{Z},\ \ -y\Big), \label{eq:urow}\\
  \text{$v$-row:}\quad
  &\frac{f}{Z}\Big(\,0,\ \ 1,\ \ \frac{y}{Z},\ \ -\frac{y^2}{Z},\ \
     \frac{xy}{Z},\ \ -x\Big). \label{eq:vrow}
\end{align}

\subsection{The normal matrix is block diagonal}

$\Jmat^{\top}\Jmat$ is the sum over the four corners of the outer products of
these rows. The square is symmetric, so sums of odd powers vanish:
\begin{equation}
  \sum_i x = \sum_i y = \sum_i xy = \sum_i x^3 = \sum_i x^2 y
  = \sum_i x y^2 = 0,
\end{equation}
while
\begin{equation}
  \sum_i x^2 = \sum_i y^2 = 4h^2,
  \qquad
  \sum_i x^4 = \sum_i y^4 = \sum_i x^2y^2 = 4h^4 .
\end{equation}
Carrying out the sums entry by entry, almost every off-diagonal term cancels,
and what remains is block diagonal:
\begin{equation}
  \frac{\Jmat^{\top}\Jmat}{(f/Z)^2}
  =
  \underbrace{\begin{bmatrix} 4 & -4s \\ -4s & 8s^2 \end{bmatrix}}
    _{(t_x,\ \omega_y)}
  \oplus
  \underbrace{\begin{bmatrix} 4 & -4s \\ -4s & 8s^2 \end{bmatrix}}
    _{(t_y,\ \omega_x)}
  \oplus
  \underbrace{\Big[\,\tfrac{8h^2}{Z^2}\,\Big]}_{t_z}
  \oplus
  \underbrace{\big[\,8h^2\,\big]}_{\omega_z},
  \qquad
  s = \frac{h^2}{Z}.
  \label{eq:blocks}
\end{equation}
For example, the $(t_x,\omega_y)$ entry is $\sum_i 1\cdot(-x^2/Z) = -4h^2/Z =
-4s$, and the $(\omega_y,\omega_y)$ entry is $\sum_i (x^4 + x^2y^2)/Z^2 =
8h^4/Z^2 = 8s^2$. The structure is already telling: \emph{each out-of-plane
rotation is coupled to one sideways translation} ($\omega_y$ with $t_x$,
$\omega_x$ with $t_y$), while depth $t_z$ and in-plane spin $\omega_z$ decouple
from everything.

\subsection{The singular values}

The singular values of $\Jmat$ are $f/Z$ times the square roots of the
eigenvalues of the blocks. The two scalar blocks give at once
\begin{equation}
  \sigma_{t_z} = \frac{2\sqrt2\,f h}{Z^2},
  \qquad
  \sigma_{\omega_z} = \frac{2\sqrt2\,f h}{Z}.
\end{equation}
Each $2\times2$ block has trace $4 + 8s^2$ and determinant
$4\cdot8s^2 - 16s^2 = 16s^2$, so its eigenvalues are
\begin{equation}
  \lambda_\pm = \frac{(4+8s^2) \pm \sqrt{(4+8s^2)^2 - 64s^2}}{2}.
\end{equation}
The marker is small compared with its distance, $s = h^2/Z \ll 1$ (here
$s = 0.005$~m), so to leading order $\lambda_+ \approx 4$ and
$\lambda_- \approx \det/\operatorname{tr} \approx 16s^2/4 = 4s^2$. Hence
\begin{equation}
  \sigma_{\max} \approx \frac{2f}{Z}
  \ \ (\text{twice}),
  \qquad
  \boxed{\;\sigma_{\min} \approx \frac{2 f h^2}{Z^2}
  \ \ (\text{twice})\;}
  \label{eq:sigmas}
\end{equation}
and the condition number is
\begin{equation}
  \kappa_{\mathrm{SI}} \approx \frac{2f/Z}{2fh^2/Z^2} = \frac{Z}{h^2}
  \quad(\text{units of } 1/\mathrm{m}).
  \label{eq:kappa}
\end{equation}
\paragraph{In the adopted units.} By \eqref{eq:Jn} the three translation columns
are multiplied by $Z$, which multiplies the translation singular values by $Z$
and leaves the rotation ones alone:
\begin{equation}
  \sigma_{t_x} = \sigma_{t_y} \approx 2f,
  \qquad
  \sigma_{t_z} = \sigma_{\omega_z} = \frac{2\sqrt2\,f h}{Z},
  \qquad
  \sigma_{\min} \approx \frac{2 f h^2}{Z^2},
  \qquad
  \boxed{\;\kappa \approx \Big(\frac{Z}{h}\Big)^2\;}
  \label{eq:sigmasn}
\end{equation}
For $\sigma_{\min}$ this is not obvious from the block alone; redoing the
$2\times2$ eigenvalue computation with $t_x$ measured in units of $Z$ gives the
same leading-order value, because the tilt direction is almost pure rotation.
The dimensionless $\kappa$ depends only on $h/Z$, the marker's apparent
half-size: \emph{the conditioning degrades as the inverse square of how large
the marker looks.} Note also that depth and in-plane spin become exactly equally
observable head-on, which Section~\ref{sec:pairs} explains.

\subsection{The weak directions, and why they are weak}

The eigenvector of $\lambda_-$ in the $(t_x,\omega_y)$ block solves
$(4-\lambda_-)\,\delta t_x - 4s\,\delta\omega_y = 0$, so
$\delta t_x \approx s\,\delta\omega_y$:
\begin{equation}
  \vvec_{5,6} \;\propto\; (s,0,0,0,1,0)^{\top}
  \quad\text{and}\quad
  (0,s,0,1,0,0)^{\top}.
  \label{eq:weakdirs}
\end{equation}
The two least observable directions are the two \emph{out-of-plane tilts},
each accompanied by a small sideways translation of $s = h^2/Z$ metres per
radian; in the adopted units that coefficient becomes $s/Z = (h/Z)^2$.

Rows \eqref{eq:urow}--\eqref{eq:vrow} show why. A tilt $\delta\omega_y$ moves
corner $i$ by
\begin{equation}
  \delta u_i = -\frac{f}{Z}\frac{x^2}{Z}\,\delta\omega_y
             = -\frac{f h^2}{Z^2}\,\delta\omega_y ,
  \qquad
  \delta v_i = \frac{f}{Z}\frac{xy}{Z}\,\delta\omega_y
             = \pm\frac{f h^2}{Z^2}\,\delta\omega_y ,
\end{equation}
because $x^2 = h^2$ at every corner. The $u$-motion is the \emph{same for all
four corners}: a uniform sideways shift, which a translation
$\delta t_x = s\,\delta\omega_y$ reproduces exactly and therefore cancels. What
is left is the $v$-motion, whose sign follows $xy$: two opposite corners move
up and the other two down, a small \emph{keystone} distortion of the square.
That keystone is the only first-order evidence of the tilt, and its size is
$f h^2/Z^2$ per radian per corner, which over four corners gives
$\sigma_{\min} = 2fh^2/Z^2$ to leading order in $s$.

Physically: head-on, tilting the marker changes its image only through
perspective, the fact that one edge comes slightly closer than the other. That
effect is of order $h/Z$ \emph{smaller} than the effect of moving the marker,
hence the factor $h^2/Z^2$ against $1/Z$.

\subsection{Check against the implementation}

With $f = 800$~px, $h = 0.05$~m, $Z = 0.5$~m, the formulas and the code agree:

\begin{center}
\begin{tabular}{@{}lrrrr@{}}
\toprule
 & \multicolumn{2}{c}{SI units ($\Jmat$)}
 & \multicolumn{2}{c}{adopted units ($\Jmat_{\mathrm{n}}$)} \\
Quantity & closed form & code & closed form & code \\
\midrule
$\sigma_{t_x} = \sigma_{t_y}$ & $3200.0$ & $3200.04$ & $1600.0$ & $1600.08$ \\
$\sigma_{t_z}$ & $452.548$ & $452.548$ & $226.274$ & $226.274$ \\
$\sigma_{\omega_z}$ & $226.274$ & $226.274$ & $226.274$ & $226.274$ \\
$\sigma_{\min}$ (twice) & $16.000$ & $15.9998$ & $16.000$ & $15.9992$ \\
$\kappa$ & $Z/h^2 = 200$ & $200.005$ & $(Z/h)^2 = 100$ & $100.010$ \\
coupling in $\vvec_{5,6}$ & $0.005$ & $0.005$ & $0.010$ & $0.010$ \\
\bottomrule
\end{tabular}
\end{center}

\noindent
The small differences in $\sigma_{\max}$ and $\kappa$ are the neglected
$O(s^2)$ terms; the exact eigenvalues reproduce the code to every digit.

% ============================================================
\section{Along the camera arc}
\label{sec:sweep}
% ============================================================

For oblique views the algebra no longer closes, and the numbers come from
\texttt{run\_conditioning\_sweep} in \texttt{experiment2.py}, which evaluates
$\Jmat_{\mathrm{n}}$ at the true pose of each viewpoint ($Z = 0.5$~m, azimuth
$0^{\circ}$, one sample per degree); they are also written to
\texttt{outputs/exp2\_conditioning.csv} and plotted in
\texttt{outputs/exp2\_conditioning.png}. Singular values are in pixels. Along
this arc $Z$ is constant, so $\kappa_{\mathrm{SI}} = \kappa/Z = 2\kappa$.

\begin{center}
\small
\begin{tabular}{@{}rrrrrrrrl@{}}
\toprule
$\theta_v$ & $\sigma_1$ & $\sigma_2$ & $\sigma_3$ & $\sigma_4$ & $\sigma_5$
  & $\sigma_6$ & $\kappa$ & weakest $\vvec_6$ \\
\midrule
$0^{\circ}$  & 1600.1 & 1600.1 & 226.3 & 226.3 & 16.0  & 16.0  & 100.0
  & $\omega_x + 0.010\,t_y$ \\
$5^{\circ}$  & 1600.3 & 1600.3 & 226.1 & 226.1 & 18.9  & 18.5  & 86.4
  & $\omega_x - 0.04\,\omega_z$ \\
$15^{\circ}$ & 1601.7 & 1601.7 & 224.6 & 224.5 & 34.1  & 32.3  & 49.6
  & $\omega_x - 0.13\,\omega_z$ \\
$30^{\circ}$ & 1606.2 & 1606.1 & 219.5 & 219.1 & 61.8  & 58.2  & 27.6
  & $\omega_x - 0.26\,\omega_z$ \\
$45^{\circ}$ & 1612.3 & 1612.1 & 210.8 & 210.0 & 89.7  & 85.1  & 18.9
  & $\omega_x - 0.38\,\omega_z$ \\
$60^{\circ}$ & 1618.4 & 1618.2 & 198.7 & 197.0 & 116.4 & 111.6 & 14.5
  & $\omega_x - 0.51\,\omega_z$ \\
$75^{\circ}$ & 1622.9 & 1622.7 & 183.2 & 180.3 & 141.0 & 136.8 & 11.9
  & $\omega_x - 0.65\,\omega_z$ \\
$89^{\circ}$ & 1624.5 & 1624.3 & 168.3 & 162.8 & 159.0 & 157.2 & 10.3
  & $\omega_z - 0.28\,\omega_x$ \\
\bottomrule
\end{tabular}
\end{center}

\subsection{Which motion each singular value belongs to}
\label{sec:pairs}

The singular vectors attach every singular value to one clear motion of the
marker, and the geometry explains why they come in pairs.
\begin{itemize}[leftmargin=*]
  \item \textbf{$\sigma_1,\sigma_2$ --- sideways translations $t_x,t_y$.}
        Sliding the marker sideways moves all four corners together. They are
        equal because a square looks the same after a $90^{\circ}$ turn, so
        sliding along $x$ or along $y$ is the same situation; along the arc the
        tilt breaks this symmetry only by $0.02\%$. They stay nearly constant
        because they depend on $f$ and $Z$ alone, both fixed on the arc; the
        slight rise comes from the near edge of a tilted marker moving closer.
  \item \textbf{$\sigma_3,\sigma_4$ --- depth $t_z$ and in-plane spin
        $\omega_z$.} Moving the marker closer by one unit of $Z$ zooms the image
        and moves each corner, at distance $r = h\sqrt2$ from the centre,
        outwards by $f r/Z$ pixels; spinning it by one radian moves each corner
        along a circle by the same $f r/Z$. Same corner motion, same singular
        value: $2\sqrt2\,fh/Z$ head-on. Both shrink as the view becomes oblique,
        because a tilted marker looks narrower (foreshortened by
        $\cos\theta_v$), and they drift apart by up to $3\%$ at grazing angles.
        In SI units the two would differ by a factor $Z$, and their order would
        depend on the unit of length.
  \item \textbf{$\sigma_5,\sigma_6$ --- the two out-of-plane tilts.} Equal
        head-on by the same $90^{\circ}$ symmetry, and tiny because a tilt then
        shows only as a keystone (Section~\ref{sec:frontal}). The arc tilts the
        marker about its own $Y$ axis, which breaks the symmetry: rotating
        further about that axis ($\omega_y$) and about the other ($\omega_x$) no
        longer look alike, and the two values separate.
\end{itemize}

\paragraph{Observations.}
\begin{itemize}[leftmargin=*]
  \item \textbf{The problem is worst conditioned head-on.} $\sigma_{\min}$
        rises tenfold from $16$ to $157$ across the sweep, and $\kappa$ falls
        from $100$ to $10$. The fronto-parallel view, often assumed to be the
        easy case, is the least observable one.
  \item \textbf{The weak directions are always rotations.} $\sigma_5$ and
        $\sigma_6$ belong to out-of-plane rotations at every angle. The
        second-weakest direction stays close to $\omega_y$, the axis of the
        viewing-angle tilt itself; the weakest starts as $\omega_x$ and
        progressively mixes with the in-plane spin $\omega_z$.
  \item \textbf{Translation is always well observed.} $\sigma_1$ and $\sigma_2$,
        the two sideways translations, stay near $2f$ throughout.
  \item \textbf{Why obliquity helps.} Head-on, an out-of-plane rotation moves
        the corners along the line of sight, and shows up only through the
        keystone, of relative size $h/Z$. Once the plane is tilted by
        $\theta_v$, the same rotation also moves each corner
        \emph{perpendicular} to the line of sight, by an amount proportional to
        $\sin\theta_v$ and with opposite signs on opposite edges. No translation
        can reproduce that, so the evidence becomes first order in $h/Z$ and
        the weakest singular value grows with the tilt.
\end{itemize}

% ============================================================
\section{Connection with Experiment 1}
\label{sec:link}
% ============================================================

The two experiments look at the same geometry from two sides. Experiment~1
showed that the two IPPE candidates differ only by the sign of $\bvec$, the line
of sight projected onto the marker plane, of length $\sin\theta_v$; the image
reveals how much the marker is tilted but not towards which side. This note
shows that, near head-on, the out-of-plane tilts are also the least observable
directions of the least-squares problem, and that their only first-order
signature is a perspective keystone of relative size $h/Z$. Both statements say
that a planar marker seen head-on constrains its own tilt weakly, and both
predict that small or distant markers make things worse.

% ============================================================
\section{What to expect from the Monte Carlo}
\label{sec:mc}
% ============================================================

Equation \eqref{eq:covsvd} turns the table into predictions. With
$\sigma_{\mathrm{px}} = 0.5$~px, the standard deviation along the weakest
direction is $\sigma_{\mathrm{px}}/\sigma_{\min}$:
\begin{center}
\begin{tabular}{@{}rrr@{}}
\toprule
$\theta_v$ & $\sigma_{\mathrm{px}}/\sigma_{\min}$ (rad) & in degrees \\
\midrule
$0^{\circ}$  & $0.0313$ & $1.79^{\circ}$ \\
$30^{\circ}$ & $0.0086$ & $0.49^{\circ}$ \\
$60^{\circ}$ & $0.0045$ & $0.26^{\circ}$ \\
\bottomrule
\end{tabular}
\end{center}
while sideways translations scatter by only $0.5/3200 \approx 0.16$~mm (these
use $\Jmat$ in SI units, as a covariance in metres and radians must; the tilt
values are the same in either convention). Two caveats follow from the rest of
this note.
\begin{enumerate}[leftmargin=*]
  \item The prediction holds for the minimizer of $E$, so the Monte Carlo must
        refine each estimate by minimizing $E$ (Gauss--Newton with $\Jmat$,
        started from IPPE), not use raw IPPE.
  \item It is a first-order, single-minimum statement. Near head-on the mirror
        candidate is almost as good as the true one (Experiment~1), so noise can
        send the estimate to the wrong minimum. There the scatter is no longer
        Gaussian around the truth, and a disagreement with \eqref{eq:covsvd} is
        expected rather than a failure.
\end{enumerate}

% ============================================================
\section{Summary}
% ============================================================

\begin{enumerate}[leftmargin=*]
  \item Residuals $\evec_i = \pi(\Kmat(\Rmat\Xvec_i+\tvec)) - \pvec_i$, stacked
        into $\evec\in\Rset^8$; cost $E = \|\evec\|^2$.
  \item Local parameters $\dth = (\dtv,\dom)$ with
        $\Rmat \mapsto \Rmat\exp(\skewm{\dom})$.
  \item $\Jmat_i = \Dmat_i\,[\,\Idm \mid -\Rmat\skewm{\Xvec_i}\,]$, with
        $\Dmat_i$ from \eqref{eq:D}; stacked, $\Jmat\in\Rset^{8\times6}$.
  \item $\Jmat = \Umat\Sigmat\Vmat^{\top}$; $\sigma_k$ is how visible the pose
        direction $\vvec_k$ is; $\sigma_{\min}$ and $\vvec_6$ identify the least
        observable direction; $\kappa = \sigma_1/\sigma_6$, computed on the
        dimensionless Jacobian $\Jmat_{\mathrm{n}} = \Jmat\Smat$ (translation in
        units of $Z$) so that it does not depend on the choice of units.
  \item Under pixel noise $\sigma_{\mathrm{px}}$, the least-squares estimate
        scatters along $\vvec_k$ with standard deviation
        $\sigma_{\mathrm{px}}/\sigma_k$.
  \item Head-on, in closed form: $\sigma_{\min} \approx 2fh^2/Z^2$ along the
        out-of-plane tilts, and $\kappa \approx (Z/h)^2$, the inverse square of
        the marker's apparent half-size ($\kappa_{\mathrm{SI}} \approx Z/h^2$).
  \item Along the arc: $\sigma_{\min}$ grows from $16$ to $157$ and $\kappa$ falls
        from $100$ to $10$; the least observable direction is always an
        out-of-plane rotation.
\end{enumerate}

\end{document}
